\chapter{分离变量法}
回忆一些典型ODE方程的基本解法
\begin{method}{二阶常系数齐次方程}{}
	形如
	\begin{equation}
		y''+ay'+by=0.
	\end{equation}
	其特征方程$\lambda^2+a\lambda+b=0$的解$\lambda_{1,2}$
	\begin{compactenum}[I.]
		\item $\Delta>0,$\quad$y=A\e^{\lambda_1x}+B\e^{\lambda_2x};$
		\item $\Delta=0,$\quad$\lambda_{1,2}=\lambda,$\quad$y=(A+Bx)\e^{\lambda x};$
		\item $\Delta<0,$\quad$\lambda_{1,2}=\alpha\pm\i\beta$,\quad$y=(A\cos\beta x+B\sin\beta x)\e^{\alpha x}$。
	\end{compactenum}
\end{method}
\begin{method}{二阶非齐次方程}{}
	形如
	\begin{equation}
		y''+a(x)y'+b(x)y=f(x).
	\end{equation}
	二阶非齐次方程
	根据叠加定理，其解$y=y_0+y_\mathrm s;$其中$y_0$为$f(x)\equiv 0$时的齐次通解，$y_\mathrm s$为特解。
	
	若$y_1,y_2$为对应的齐次方程线性无关解，则可求出一特解
	\begin{align}\label{General to Particular Solution}
		y_\mathrm s(x)=\int_{x_0}^x\frac{y_1(\xi)y_2(x)-y_1(x)y_2(\xi)}{y_1(\xi)y_2'(\xi)-y_1'(\xi)y_2(\xi)}f(\xi)\d \xi,\quad\forall x_0.
	\end{align}
	特解满足齐次边界条件
	\[
		y_\mathrm s(x_0)=0,\quad y'_\mathrm s(x_0)=0.
	\]
\end{method}
\begin{method}{Euler方程}{}
	形如
	\begin{equation}
		x^2y''+axy'+by=f(x).
	\end{equation}
	令$x=\e^t,\;u(t)=y(\e^t)=y(x)$方程变为
	\begin{align}
		u''+(a-1)u'+bu=f(\e^t).
	\end{align}
\end{method}

\section{特征值和特征函数}
在讨论分离变量前，先引入Fourier级数的概念。
\begin{definition}{Fourier级数}{Fourier Series}
	$f(x)$周期为$2\ell$，则$f(x)$的Fourier级数
	\begin{align}
		f(x)\sim\frac{a_0}2+\sum_{n=1}^\infty\Bigp{a_n\cos\frac{n\pi x}\ell+b_n\sin\frac{n\pi x}\ell}.
	\end{align}
	其中Fourier系数
	\begin{align*}
		a_n & =\frac1\ell\int_{-\ell}^\ell f(x)\cos\frac{n\pi x}\ell\d x; \\
		b_n & =\frac1\ell\int_{-\ell}^\ell f(x)\sin\frac{n\pi x}\ell\d x.
	\end{align*}
\end{definition}
\begin{theorem}{Parseval等式}{Parserval Equality}
	若$f\in L^2\bigb{-\ell,\ell}$即平方可积，则
	\[
		\frac1\ell\int_{-\ell}^\ell \abs{f(x)}^2\d x=\frac{\abs{a_0}^2}2+\sum_{n=1}^\infty(\abs{a_n}^2+\abs{b_n}^2).
	\]
\end{theorem}
\begin{theorem}{Dirichlet收敛定理}{Dirichlet Convergence Theorem}
	若
	\begin{compactenum}
		\item $f,f'$连续或分段连续，至多有有限个第一类间断点；%（可去或跳跃）
		\item $f$至多有有限个极值点；
	\end{compactenum}
	则
	\begin{align*}
		f(x)
		\sim\begin{cases}
			\bigb{f(x^+)+f(x^-)}/2,        & x\in(-\ell,\ell) \\
			\bigb{f(-\ell^+)+f(\ell^-)}/2, & x=\pm\ell
		\end{cases}
	\end{align*}
	在$f$连续点处，$f(x)$与其Fourier级数相等。
\end{theorem}
\begin{example}{有界弦问题}{}
	有界弦
	\begin{align*}
		\begin{cases}
			u_{tt}=a^2u_{xx},\quad x\in(0,L),\;t>0 \\
			u(0,t)=u(L,t)=0,                 \\
			u(x,0)=\varphi(x)\quad u_t(x,0)=\psi(x).
		\end{cases}
	\end{align*}
	假设变量是可以分离的，即$u(x,t)=X(x)T(t),$则引入\textbf{特征值}$\lambda$
	\[
		\frac{X''}X=\frac{T''}{a^2T}=-\lambda.
	\]
	要使$X''+\lambda X=0$且$X(0)=X(L)=0$仅有
	\[
		\lambda_n=\Bigp{\frac{n\pi}L}^2,\quad X_n(x)=\sin\frac{n\pi x}L,\quad n=1,2,\ldots.
	\]
	进而$T''+a^2\lambda T=0$
	\[
		T_n(t)=C_n\cos\frac{n\pi at}L+D_n\sin\frac{n\pi at}L.
	\]
	因此原方程应有\textit{形式解}
	\begin{align}
		u(x,t)=\sum_{n=1}^\infty\Bigb{C_n\cos\Bigp{\frac{n\pi at}L}+D_n\sin\Bigp{\frac{n\pi at}L}}\sin\Bigp{\frac{n\pi x}L}.
	\end{align}
	$C_n,D_n$由初始条件确定
	\begin{align*}
		u(x,0)&=\sum_{n=1}^\infty C_n\sin\Bigp{\frac{n\pi x}L}=\varphi(x),            \\
		u_t(x,0)&=\sum_{n=1}^\infty D_n\frac{n\pi a}L\sin\Bigp{\frac{n\pi x}L}=\psi(x).
	\end{align*}
	恰好对应Fourier正弦系数$\varphi^{\sine}_n,\psi^{\sine}_n$
	\begin{align*}
		C_n & =\frac2L\int_0^L\varphi(x)\sin\Bigp{\frac{n\pi x}L}\d x=:\varphi_n^{\sine},                  \\
		D_n & =\frac2{n\pi a}\int_0^L\psi(x)\sin\Bigp{\frac{n\pi x}L}\d x=:\frac{L}{n\pi a}\psi_n^{\sine}.
	\end{align*}
	
	将形式级数当作古典解时，还须检查逐项微分的收敛性及初边值相容性；仅有平方可积初值时，应在相应的弱解意义下理解展开。
\end{example}
\begin{example}{圆域Laplace方程}{}
	圆域上的Laplace方程，采用极坐标
	\begin{align*}
		\begin{cases}
			\laplacian u=u_{rr}+\frac1ru_r+\frac1{r^2}u_{\theta\theta}=0,&r<r_0\\
			u(r_0,\theta)=\varphi(\theta)
		\end{cases}
	\end{align*}
	还应包括自然条件，即
	\[
		\limsup_{r\to0^+}\abs{u(r,\theta)}<\infty.
	\]
	和周期边界条件
	\[
		u(r,0)=u(r,2\pi),\quad u_\theta(r,0)=u_\theta(r,2\pi).
	\]
	
	分离变量$u(r,\theta)=R(r)\varTheta(\theta)$
	\begin{align*}
		\begin{cases}
			r^2R''+rR'-\lambda R=0\\
			\varTheta''+\lambda\varTheta=0.
		\end{cases}
	\end{align*}
	由周期条件，特征值$\lambda=n^2,\;n=0,1,2,\ldots$
	\[
		\varTheta_n(\theta)=a_n\cos n\theta+b_n\sin n\theta.
	\]
	再解$R$的Euler方程，可得
	\[
		R_n(r)=
		\begin{cases}
			c_0+d_0\ln r,&n=0\\
			c_nr^n+d_nr^{-n},&n\geqslant 1
		\end{cases}
	\]
	再由自然条件$\abs{R(0)}<\infty$，可知$d_i\equiv 0$，进而
	\[
		u(r,\theta)=\frac{a_0}2+\sum_{n=1}^\infty r^n(a_n\cos n\theta+b_n\sin n\theta),
	\]
	结合边界条件$u(r_0,\theta)=\varphi(\theta)$，求得系数
	\[
		a_n=\frac1{\pi r_0^n}\int_0^{2\pi}\varphi(\vartheta)\cos n\vartheta\d\vartheta,\quad b_n=\frac1{\pi r_0^n}\int_0^{2\pi}\varphi(\vartheta)\sin n\vartheta\d\vartheta.
	\]
	因此
	\begin{align*}
		u(r,\theta)&=\frac1{2\pi}\int_0^{2\pi}\varphi(\vartheta)\d\vartheta+\\
		&\qquad\frac1\pi\sum_{n=1}^\infty\Bigp{\frac r{r_0}}^n\int_0^{2\pi}\varphi(\vartheta)(\cos n\vartheta\cos n\theta+\sin n\vartheta\sin n\theta)\d\vartheta\\
		&=\frac1{2\pi}\int_0^{2\pi}\varphi(\vartheta)\biggb{1+2\sum_{n=1}^\infty\Bigp{\frac r{r_0}}^n\cos n(\theta-\vartheta)}\d\vartheta
	\end{align*}
	又由Euler公式
	\begin{align*}
	\sum_{n=0}^\infty k^n\cos n\theta=\sum_{n=0}^\infty\Re k^n\e^{\i n\theta}=\Re\frac1{1-k\e^{\i\theta}}=\frac{1-k\cos\theta}{1+k^2-2k\cos\theta}.
	\end{align*}
	故得到圆域内的Poisson公式
	\begin{align}
		u(r,\theta)&=\frac1{2\pi}\int_0^{2\pi}\frac{(r_0^2-r^2)\varphi(\vartheta)}{r_0^2+r^2-2r_0r\cos(\theta-\vartheta)}\d\vartheta
	\end{align}
	积分内式子称为Poisson核。
\end{example}
\begin{example}{圆域Neumann问题}{}
	圆域上的Neumann问题
	\begin{align*}
		\begin{cases}
			\laplacian u=0,\quad r<R\\
			\biggeval{\pv un}_{r=R}=\varphi(\theta)
		\end{cases}
	\end{align*}
	由~\textcolor[rgb]{0,.8,0}{Green公式}
	\begin{align*}
		{\color[rgb]{0,.8,0}{\int_{D_R}\laplacian u\d x\nd y=\oint_{r=R}\pv un\d\ell}}=R\int_0^{2\pi}\varphi(\theta)\d\theta=0.
	\end{align*}
	剩下的证明留给读者
\end{example}

\section{Sturm--Liouville定理}
\begin{example}{有限杆热传导}{}
	设杆温度$u=u(x,t)$，则定解问题
	\begin{align*}
		\begin{cases}
			u_t=a^2u_{xx},\quad x\in(0,L),\;t>0,\\
			u(0,t)=0,\quad u_x(L,t)+hu(L,t)=0,\\
			u(x,0)=\varphi(x).
		\end{cases}
	\end{align*}
	其中$h>0$是表面换热系数与热导率之比，量纲为长度的倒数，环境温度取为零。
	
	分离常数$u(x,t)=X(x)T(t),$
	\begin{align*}
		\begin{cases}
			X''+\beta^2X=0, \\
			X(0)=0,\quad X'(L)+hX(L)=0.
		\end{cases}\implies\beta+h\tan\beta L=0.
	\end{align*}
	$\implies X_n(x)=\sin\beta_nx,\quad T_n(t)=C_n\e^{-\beta_n^2a^2t}$。
	
	可以验证$X_n$是正交的，即
	\[
		\inp{X_m}{X_n}=\int_0^LX_m(x)X_n(x)\d x=0,\quad m\neq n.
	\]
	形式解
	\[
		u(x,t)=\sum_{n=1}^\infty C_n\e^{-\beta_n^2a^2t}\sin\beta_nx.
	\]
	
	再根据初始条件
	\[
		u(x,0)=\sum_{n=1}^\infty C_n\sin\beta_nx=\varphi(x).
	\]
	在后面可以看到，这是广义Fourier级数，其系数
	\[
		C_n=\frac{\inp{\varphi}{\sin\beta_nx}}{\inp{\sin\beta_nx}{\sin\beta_nx}}=\frac{\int_0^L\varphi(x)\sin\beta_nx\d x}{\int_0^L\sin^2\beta_nx\d x}.
	\]
\end{example}
\begin{definition}{加权内积}{Weight Inner Product}
	定义在$\bigb{a,b}$上的实函数$f,g$的$\rho$\;-\;加权内积
	\[
		\inp fg_\rho:=\int_a^bf(x)g(x)\rho(x)\d x.
	\]
	其中权函数$\rho\geqslant 0$分段连续且零点孤立。
	% 若$f,g$正交，则$\inp fg_\rho=0$。
	%定义范数$\norm f_\rho:=\inp ff_\rho^{1/2}$。及

	加权平方可积函数空间把几乎处处相等的函数视为同一个元素，记作
	\[
		L^2_\rho\bigb{a,b}:=\set f{\norm f_\rho<\infty}.
	\]
\end{definition}
\begin{theorem}{广义Fourier级数}{Generalized Fourier Series}
	若$f_1,f_2,\ldots$在$L^2_\rho$中完备且加权正交，则$\forall f\in L^2_\rho$有广义Fourier级数展开
	\[
		f(x)\sim\sum_{i=1}^\infty a_if_i(x),\quad a_i=\frac{\inp f{f_i}_\rho}{\norm{f_i}^2_\rho}.
	\]
	且有Parseval等式
	\[
		\norm f_\rho^2=\sum_{i=1}^\infty\bigginp f{\frac{f_i}{\norm{f_i}_\rho}}_\rho^2.
	\]
\end{theorem}
\begin{definition}{Sturm--Liouville方程}{Sturm-Liouville Equation}
	定义域$[a,b]$，常见一维特征值问题都可以化为
	\[
		\bigb{k(x)f'(x)}'-q(x)f(x)+\lambda\rho(x)f(x)=0.
	\]
	其中$\lambda$为参数，$k,q,\rho$为实函数。
	记
	\[
		\SLop[f]:=-\frac{(kf')'-qf}\rho.
	\]
	则
	\[
		\inp{\SLop[f]}g_\rho-\inp f{\SLop[g]}_\rho=\bigb{k(fg'-f'g)}_a^b.
	\]
	好的边界条件可以使上式等于0。
\end{definition}
\begin{definition}{正则Sturm--Liouville问题}{}
    正则Sturm--Liouville问题
    \[
        \SLop[f](x)=\lambda f(x).
    \]
    
    正则条件：保证Sturm--Liouville方程在$\bigb{a,b}$上没有奇点。
    \[
    	k\in\cont^2\bigb{a,b},\;q,\rho\in\cont\bigb{a,b};\;k,\rho>0.
    \]
    
    
    好的边界条件：
    \begin{center}
    	可分：$\begin{cases}
    			c_1f(a)+c_2f'(a)=0, \\
    			d_1f(b)+d_2f'(b)=0.
    		\end{cases}$\quad 或 \quad 周期：
    	$\begin{cases}
    			f(a)=f(b) \\
			k(a)f'(a)=k(b)f'(b)
    		\end{cases}$
    \end{center}
    可分边界的实系数对$(c_1,c_2),(d_1,d_2)$均不同时为零。上述条件可使$\SLop$是对称算子，即
    \[
        \inp{\SLop[f]}g_\rho=\inp f{\SLop[g]}_\rho.
    \]
\end{definition}
\begin{theorem}{Sturm--Liouville定理1}{Sturm-Liouville Theorem I}
	正则Sturm--Liouville问题
	\begin{compactenum}
		\item 有可数多个实特征值（以下按互异特征值编号）
		      \[
				\lambda_1<\lambda_2<\cdots<\lambda_n<\cdots,\quad\lim_{n\to\infty}\lambda_n=+\infty.
			  \]
		      % 对应特征函数
		      % \[X_1(x),X_2(x),\ldots,X_n(x),\ldots\]
		\item 不同特征值对应的特征函数加权$\rho$\;-\;正交；同一特征子空间内可选取正交基；
		\item 特征值的特征子空间至多2维，

		      可分边界条件下，特征子空间为1维；
		\item 特征函数构成$L^2_\rho$上完备的正交基底。
	\end{compactenum}
    \tcblower
    奇异Sturm--Liouville问题不满足上述正则条件，例如端点处系数退化或区间无界。此时必须另外确定端点条件和算子的自伴定义域，谱可能含连续部分，不能直接套用上述离散谱结论。
\end{theorem}
\begin{theorem}{Sturm--Liouville定理2}{Sturm-Liouville Theorem II}
	若正则Sturm--Liouville问题还满足$q\geqslant 0$；可分边界条件下还需满足$c_1c_2\leqslant 0,\;d_1d_2\geqslant 0$。
	\begin{compactenum}
		\item 所有特征值非负；

		      特别地，在上述可分边界条件下，存在0特征值的充要条件是$q\equiv 0$且两端为Neumann条件，此时常函数为特征函数；周期边界下，$q\equiv 0$也给出常函数零模；
		\item 周期边界下特征值可为单重或二重，不能断言所有非最低特征值均二重简并。例如常系数问题$-f''=\lambda f$在周期区间上具有正弦、余弦两种同频模式，但一般变系数问题不一定保留这一简并。
	\end{compactenum}
\end{theorem}
\begin{example}{扇形域Dirichlet问题}{}
	扇形域上的Dirichlet问题
	\begin{align*}
		\begin{cases}
			\laplacian u=0,\quad r\in(1,\e),\;\theta\in(0,\pi/2).\\
			u(1,\theta)=u(\e,\theta)=0,\\
			u(r,0)=0,\quad u(r,\pi/2)=g(r).
		\end{cases}
	\end{align*}
	化为Sturm--Liouville型
	\[
		(rR')'+\frac{\lambda}rR=0,
	\]
	对应\(k=r,\;q=0,\;\rho=\frac1r\)，因此\(\lambda>0\)，进而
	\begin{gather*}
		\lambda_n=(n\pi)^2,\quad R_n(r)=\sin(n\pi\ln r),\quad n=1,2,\ldots
		\\
		\varTheta_n(\theta)=A_n\e^{n\pi\theta}+B_n\e^{-n\pi\theta}.
	\end{gather*}
	再结合$u$的边界条件即可解出。
\end{example}

\section{非齐次方程}

对非齐次方程，有时可靠直觉找出特解$v$，但对一般的非齐次方程，需要借助齐次化原理转化为齐次方程：

\begin{method}{齐次化原理法}{}
	设$\DiffOp_t$中关于2阶导数的系数为1
	\begin{align*}
		\begin{cases}
			\DiffOp_xu+\DiffOp_tu=f(x,t),\quad x\in(a,b),\;t>0\\
			(\alpha_1u-\beta_1u_x)_{x=a}=0,\quad(\alpha_2u+\beta_2u_x)_{x=b}=0,\\
			\eval u_{t=0}=\eval{u_t}_{t=0}=0.
		\end{cases}
	\end{align*}
	由 \pageref{thm:Homogeneity Principle} 页的\thmref{thm:Homogeneity Principle} 知，有解
	\[
		u(x,t)=\int_0^tv(x,t-\tau;\tau)\d\tau
	\]
	$v$满足
	\begin{align*}
		\begin{cases}
			\DiffOp_xv+\DiffOp_tv=0,\quad x\in(a,b),\;t>0\\
			(\alpha_1v-\beta_1v_x)_{x=a}=0,\quad(\alpha_2v+\beta_2v_x)_{x=b}=0,\\
			\eval v_{t=0}=0,\quad\eval{v_t}_{t=0}=f(x,\tau).
		\end{cases}
	\end{align*}
\end{method}

\begin{method}{广义Fourier展开法}{}
	对于 
	\begin{align*}
		\begin{cases}
			\DiffOp_xu+\DiffOp_tu=f(x,t),\quad x\in(a,b),\;t>0\\
			(\alpha_1u-\beta_1u_x)_{x=a}=0,\quad(\alpha_2u+\beta_2u_x)_{x=b}=0,\\
			\eval u_{t=0}=\varphi(x),\quad\eval{u_t}_{t=0}=\psi(x).
		\end{cases}
	\end{align*}
	\begin{compactenum}
		\item 首先分离变量，求出$f\equiv 0$对应的特征值问题
		\[
			\DiffOp_xX=-\lambda X.
		\]
		解出特征值$\{\lambda_n\}$和本征函数$\{X_n(x)\}$，注意$n$的取值是否含0。
		\item 根据Sturm--Liouville定理判断$\{X_n\}$的完备性，展开
		\begin{gather*}
			u(x,t)=\sum T_n(t)X_n(x),\quad f(x)=\sum f_n(t)X_n(x),\\
			\varphi(x)=\sum\varphi_nX_n(x),\quad \psi(x)=\sum\psi_nX_n(x).
		\end{gather*}
		求解 
		\begin{align*}
			\begin{cases}
				\DiffOp_tT_n(t)-\lambda_nT_n(t)=f_n(t),\\
				T_n(t)=\varphi_n,\quad T_n'(t)=\psi_n.
			\end{cases}
		\end{align*}
		\item 解出$T_n(t)$，得$u(x,t)$
	\end{compactenum}
\end{method}

\begin{method}{一般地非齐次问题}{}
	对于一般地非齐次问题
	\begin{align*}
		\begin{cases}
			\DiffOp_xu+\DiffOp_tu=f(x,t),\quad x\in(a,b),\;t>0\\
			(\alpha_1u-\beta_1u_x)_{x=a}=g_1(t),\quad(\alpha_2u+\beta_2u_x)_{x=b}=g_2(t),\\
			\eval u_{t=0}=\varphi(x),\quad\eval{u_t}_{t=0}=\psi(x).
		\end{cases}
	\end{align*}
	首先将非齐次边界条件齐次化，寻找$v(x,t)$满足 
	\[
		(\alpha_1v-\beta_1v_x)_{x=a}=g_1(t),\quad(\alpha_2v+\beta_2v_x)_{x=b}=g_2(t).
	\]
	
	显然$v$不唯一，若无法直觉看出，可设为$x$的线性函数
	\[
		v(x,t)=A(t)x+B(t),
	\]
	整理得
	\[
		\begin{cases}
			(\alpha_1a-\beta_1)A(t)+\alpha_1B(t)=g_1(t),\\
			(\alpha_2b+\beta_2)A(t)+\alpha_2B(t)=g_2(t).
		\end{cases}
	\]
	便可解出$A,B$若无解还可设为$x$的二次函数。
	
	这样便可将$u$分解为$u=v+w$，而$w$的边界条件是齐次化的：
	\begin{align*}
		\begin{cases}
			\DiffOp_xw+\DiffOp_tw=f(x,t)-\DiffOp_xv-\DiffOp_tv,\quad x\in(a,b),\;t>0\\
			(\alpha_1w-\beta_1w_x)_{x=a}=0,\quad(\alpha_2w+\beta_2w_x)_{x=b}=0,\\
			\eval w_{t=0}=\varphi(x)-\eval v_{t=0},\quad\eval{w_t}_{t=0}=\psi(x)-\eval{v_t}_{t=0}.
		\end{cases}
	\end{align*}
	这在前面已解决。
\end{method}

